\medskip

On considère la fonction $f$  définie sur $]0~;~+\infty[$ par 
$f(x) = x - x \ln x,$

où ln désigne la fonction logarithme népérien.

\bigskip

\textbf{Partie A}

\medskip

\begin{enumerate}
\item On cherche la limite de $f(x)$ quand $x$ tend vers $0$.

On sait que $\ds\lim_{x\to 0 \atop x>0} x\,\ln x=0$ donc $\ds\lim_{x\to 0 \atop x>0} f(x) = 0$.

\item On cherche la limite de $f(x)$ quand $x$ tend vers $+\infty$.

$f(x)=x\left (1 - \ln x\right )$.

$\ds\lim_{x\to +\infty} \ln x = +\infty$ donc $\ds\lim_{x\to +\infty} 1- \ln x = -\infty$

Par produit de limites, on déduit : $\ds\lim_{x\to +\infty} f(x) = -\infty$

\item On admet que la fonction $f$ est dérivable sur $]0~;~+\infty[$ et on note $f'$ sa fonction dérivée.
	\begin{enumerate}
		\item Pour tout réel $x > 0$, on a :
		$f'(x) = 1 - 1\times \ln x - x\times \dfrac{1}{x} = - \ln x$.
		
		\item 
		
On détermine le signe de $f'(x)$:

\begin{center}
{\renewcommand{\arraystretch}{1.5}
\def\esp{\hspace*{3cm}}
$\begin{array}{|c |l *{4}{c} |} 
\hline
x  & 0 & \esp & 1 & \esp  & +\infty \\
\hline
\ln x &\vline\;\vline  & \pmb{-} &  \vline\hspace{-2.7pt}{0} & \pmb{+} &   \\
\hline
f'(x) &\vline\;\vline  & \pmb{+} &  \vline\hspace{-2.7pt}{0} & \pmb{-} &   \\
\hline
\end{array}$
}
\end{center}
		
$f(1) = 1 -\ln 1=1$; on établit le tableau des variations de la fonction $f$:		
		
\begin{center}
{\renewcommand{\arraystretch}{1.3}
\psset{nodesep=3pt,arrowsize=2pt 3}
\def\esp{\hspace*{2.5cm}}
\def\hauteur{0pt}
$\begin{array}{|c|l*4{c}|}
\hline
x 		& 0  & \esp & 1 & \esp & +\infty \\\hline
f'(x) 	& \vline\;\vline &   \pmb{+} & \vline\hspace{-2.7pt}0 & \pmb{-} & \\\hline
 		&  \vline\;\vline & & \Rnode{max}{1}  & & \\
f(x) &\vline\;\vline & &  & & \rule{0pt}{\hauteur} \\
 &\vline\;\vline \Rnode{min1}{~0} & & & & \Rnode{min2}{-\infty} \rule{0pt}{\hauteur}    
 \ncline{->}{min1}{max}
 \ncline{->}{max}{min2}\\ \hline
\end{array} $
}
\end{center}	
		
	\end{enumerate}
\item On résout l'équation $f(x) = x$ sur $]0~;~+\infty[$.

$f(x)=x \iff x - x\ln x = x \iff -x \ln x=0 \iff \ln x = 0 \iff x = 1$
\end{enumerate}

\bigskip

\textbf{Partie B}

\medskip

Dans cette partie, on pourra utiliser avec profit certains résultats de la partie A. 

On considère la suite $\left(u_n\right)$ définie par:

\[\left\{\begin{array}{l !{=} l}
u_0 & 0,5\\
u_{n+1} & u_n - u_n \ln u_n \text{ pour tout entier naturel } n,
\end{array}\right.\]

Ainsi, pour tout entier naturel $n$, on a : $u_{n+1} = f\left(u_n\right)$.



\begin{enumerate}
\item On rappelle que la fonction $f$ est croissante sur l'intervalle [0,5~;~1].



Soit $\mathcal{P}_n$ la propriété: $0,5 \leqslant  u_n \leqslant u_{n+1} \leqslant 1$.

\begin{list}{\textbullet}{}
\item Pour $n=0$, $u_n=u_0=0,5$ et $u_1=f(u_0)=f(0,5)\approx 0,85$.

On a donc $0,5 \leqslant u_0 \leqslant u_1 \leqslant 1$, ce qui vaut dire que $\mathcal{P}_0$ est vraie.

\item On suppose la propriété vraie au rang $n$, c'est-à-dire que $0,5 \leqslant  u_n \leqslant u_{n+1} \leqslant 1$.

La fonction $f$ est croissante sur  $[0,5~;~1]$ donc
$f(0,5) \leqslant  f(u_n) \leqslant f(u_{n+1}) \leqslant f(1)$.

Or $f(0,5)\approx 0,85 \geqslant 0,5$, $f(u_n)=u_{n+1}$, $f(u_{n+1}) = u_{n+2}$ et $f(1)=1$.

On a donc: $0,5 \leqslant u_{n+1} \leqslant u_{n+2} \leqslant 1$, donc  la propriété $\mathcal{P}_{n+1}$ est vraie.

\item La propriété $\mathcal{P}_n$ est vraie au rang 0; de plus elle est héréditaire. Donc, d'après le principe de récurrence, la propriété est vraie pour tout entier naturel $n$.
\end{list}

On a donc démontré que, pour tout entier naturel $n$, on a : $0,5 \leqslant  u_n \leqslant u_{n+1} \leqslant 1$.

\item 
	\begin{enumerate}
		\item 
Des inégalités $0,5 \leqslant  u_n \leqslant u_{n+1} \leqslant 1$, on peut déduire:

\begin{list}{\textbullet}{}
\item  $u_n \leqslant u_{n+1}$ donc la suite $(u_n)$ est croissante;
\item $u_n\leqslant 1$ donc la suite $(u_n)$ est majorée par 1.
\end{list}

D'après le théorème de la convergence monotone, on peut en déduire que la suite $(u_n)$ est convergente.
		
		
		\item On note $\ell$ la limite de la suite $\left(u_n\right)$. 
		
Pour tout $n$, on a $f(u_n)= u_{n+1}$. La fonction $f$ étant continue, la limite $\ell$ de la suite $(u_n)$ vérifie donc $f(\ell) = \ell$.

L'équation $f(x)=x$ a pour solution $x=1$ donc on peut dire que $\ell=1$.
		
	\end{enumerate}
\end{enumerate}

\bigskip

\textbf{Partie C}

\medskip

Pour un nombre réel $k$ quelconque, on considère la fonction $f_k$ définie sur $]0~;~+\infty[$ par:
$f_k(x) = kx - x \ln x.$

\begin{enumerate}
\item 

On étudie les variations de la fonction $f_k$.

$f'_k(x)=k - \ln x - x\times \dfrac{1}{x}= k-1 -\ln x$

\textbullet~~$f'_k(x) = 0 \iff k-1-\ln x = 0 \iff k-1 = \ln x \iff \e^{k-1}=x$

\textbullet~~$f'_k(x) > 0 \iff k-1-\ln x > 0 \iff k-1 > \ln x \iff \e^{k-1}>x$

Donc $f'_k(x)>0$ si $x<\e^{k-1}$, $f'_k(x)=0$ si $x=\e^{k-1}$ et $f'_k(x)<0$ si $x>\e^{k-1}$.

On en déduit que la fonction $f_k$ admet un maximum $y_k$ atteint pour $x_k=\e^{k-1}$.

\item  
$x_k = \e^k- 1$ donc 

$\begin{gathered}y_k=f_k(x_k) = kx_k - x_k \ln x_k
=k \e^{k- 1} -  \e^{k- 1}\ln \e^{k- 1}
= k\e^{k-1} - \e^{k- 1}\left (k-1\right )\\
\phantom{y_k}
=  k\e^{k-1} -  k\e^{k-1} +  \e^{k-1}
=\e^{k-1}
=x_k
\end{gathered}$


\end{enumerate}

\bigskip
